\chapter{Stylised Facts of Returns}\label{ch:rs:stylised-facts-of-returns}

On 19 October 1987 the Dow Jones Industrial Average fell 22.6\% in one session. The whole US stock market, value-weighted and in
excess of bills, fell 17.4\%: 22 times the standard deviation of its daily returns over the previous five years. Under a normal law
such a day has a probability too small to be worth writing down; in a century of daily data it is one of 103 days beyond five
standard deviations, where a normal law would expect 0.015. Real returns are not normal, and the ways in which they are not are
regular enough to have a name. This chapter states the \omterm{def:rs:stylised-facts-of-returns:fact}{stylised facts} of returns and measures them on a century of US data; it
then builds the book's synthetic equity market and checks which facts it reproduces and which it misses, because every later
chapter that uses the simulator inherits both.

\section{Heavy tails and aggregational Gaussianity}

\begin{definition}[Stylised fact]\label{def:rs:stylised-facts-of-returns:fact}
A \emph{stylised fact}\index{stylised fact} is a statistical property of returns found, qualitatively, across many assets,
markets and periods, stated without the precision that would make it specific to one of them.
\end{definition}

Cont (2001) lists eleven. The first two are the absence of linear autocorrelation of returns beyond very short intraday horizons,
and heavy tails: the tails of the return distribution decay like a power, with a tail index (One Quant Book 4, chapter 15) above
two and below five for most data sets. The daily excess return of the US market from July 1926 to July 2026 (26\,296 days, from the
Kenneth French data library) has a kurtosis of 19.1; the Hill estimator on the 500 largest losses gives a tail index of 2.90, and
on the 500 largest gains 2.58. Gopikrishnan and co-authors (1998) found an exponent of about three in forty million trade-by-trade
price changes of US stocks, the ``inverse cubic law''. \Cref{fig:rs:stylised-facts-of-returns:hist} shows the century's daily
returns against a normal law with the same standard deviation.

\begin{omfigure}
\begin{tikzpicture}
\begin{semilogyaxis}[width=10.5cm,height=5.4cm,xlabel={daily excess return in standard deviations},ylabel={frequency (log scale)},
  tick label style={font=\small},label style={font=\small},xmin=-10.5,xmax=10.5,ymin=1e-5,ymax=1,grid=major,grid style={gray!25},
  legend columns=2,legend style={font=\footnotesize,at={(0.5,-0.32)},anchor=north,draw=none,fill=none,
  /tikz/every even column/.append style={column sep=8pt}},table/col sep=comma,unbounded coords=jump]
\addplot[omDef,only marks,mark=*,mark size=1.3pt] table[x=z,y=freq]{figdata/research/05-stylised-facts-of-returns/hist.csv};
\addplot[omThm,thick,dashed] table[x=z,y=normal]{figdata/research/05-stylised-facts-of-returns/hist.csv};
\legend{US market 1926--2026,normal law}
\end{semilogyaxis}
\end{tikzpicture}

\omcaption{Frequency of the US market's daily excess returns in bins of half a standard deviation, July 1926 to July 2026, against
the normal law with the same mean and standard deviation. The normal law gives the bins beyond five standard deviations a combined
expected count of 0.015; the century has 103 such days. Data: Kenneth R.~French data library (derived statistics).}\label{fig:rs:stylised-facts-of-returns:hist}
\end{omfigure}

\begin{definition}[Aggregational Gaussianity]\label{def:rs:stylised-facts-of-returns:aggregational}
\emph{Aggregational Gaussianity}\index{aggregational Gaussianity} is the tendency of returns summed over longer horizons to look
more normal: their kurtosis falls towards three as the horizon grows, so that the shape of the distribution depends on the
horizon.
\end{definition}

The fall is slow (\cref{fig:rs:stylised-facts-of-returns:agg}). The kurtosis of the century's log returns is 20.0 at one day,
11.5 at a week, 9.5 at a month and 7.3 at a quarter. Independent draws would bring a kurtosis of 20 down to $3 + 17/21 = 3.8$ in a
month of 21 days; the returns are uncorrelated but not independent, and volatility clustering keeps the monthly tails heavy.

\begin{omfigure}
\begin{tikzpicture}
\begin{semilogxaxis}[width=10cm,height=5cm,xlabel={horizon (trading days, log scale)},ylabel={kurtosis of log returns},
  tick label style={font=\small},label style={font=\small},xmin=0.8,xmax=80,ymin=0,ymax=22,grid=major,grid style={gray!25},
  xtick={1,5,10,21,63},xticklabels={1,5,10,21,63},legend columns=3,legend style={font=\footnotesize,at={(0.5,-0.32)},
  anchor=north,draw=none,fill=none,/tikz/every even column/.append style={column sep=8pt}},table/col sep=comma]
\addplot[omDef,very thick,mark=*] table[x=horizon,y=french]{figdata/research/05-stylised-facts-of-returns/agg.csv};
\addplot[omThm,very thick,dashed,mark=square*] table[x=horizon,y=synth]{figdata/research/05-stylised-facts-of-returns/agg.csv};
\addplot[black!60,thin,sharp plot] coordinates {(0.8,3) (80,3)};
\legend{US market 1926--2026,synthetic market factor,normal}
\end{semilogxaxis}
\end{tikzpicture}

\omcaption{Kurtosis of log returns summed over non-overlapping horizons. The US market's falls from 20.0 at one day to 9.5 at a
month and 7.3 at a quarter; the book's synthetic market factor (ten years, GJR-GARCH with Student-$t$ shocks) falls from 18.3 to 3.8
at a month. Data: Kenneth R.~French data library (derived statistics); \texttt{firm.synthmkt}, seed 1.}\label{fig:rs:stylised-facts-of-returns:agg}
\end{omfigure}

\section{Volatility clustering and long memory in size}

Returns are nearly uncorrelated: the century's first-order autocorrelation is 0.046. Their sizes are not. Large moves follow
large moves (volatility clustering, Book 4, chapter 18), and the dependence lasts: the autocorrelation of absolute returns is 0.30
at one day, 0.22 at twenty days and 0.13 at a hundred (\cref{fig:rs:stylised-facts-of-returns:acf}). Cont describes the decay as
roughly a power of the lag with an exponent between 0.2 and 0.4, the long memory of Book 4, chapter 17.

\begin{definition}[Taylor effect]\label{def:rs:stylised-facts-of-returns:taylor}
The \emph{Taylor effect}\index{Taylor effect} is the finding that the autocorrelations of absolute returns exceed those of their
powers $|r|^\theta$ for $\theta \ne 1$, in particular those of squared returns (Granger and Ding, 1996, after Taylor, 1986).
\end{definition}

The century shows it at every lag: 0.30 against 0.26 at one day, 0.22 against 0.11 at twenty, 0.13 against 0.04 at a hundred.
Squares give the largest moves too much weight, and the largest moves are the least persistent part of volatility. The practical
consequence is modest but real: models and diagnostics of volatility persistence built on absolute returns see more of it than
those built on squares.

\begin{omfigure}
\begin{tikzpicture}
\begin{semilogxaxis}[width=10.5cm,height=5.2cm,xlabel={lag (trading days, log scale)},ylabel={autocorrelation},
  tick label style={font=\small},label style={font=\small},xmin=1,xmax=250,ymin=-0.05,ymax=0.35,grid=major,grid style={gray!25},
  legend columns=2,legend cell align=left,legend style={font=\footnotesize,at={(0.5,-0.32)},anchor=north,draw=none,fill=none,
  /tikz/every even column/.append style={column sep=8pt}},table/col sep=comma]
\addplot[omDef,very thick,mark=*,mark size=1.2pt] table[x=lag,y=abs]{figdata/research/05-stylised-facts-of-returns/acf.csv};
\addplot[omDef!60,thick,mark=o,mark size=1.2pt] table[x=lag,y=sq]{figdata/research/05-stylised-facts-of-returns/acf.csv};
\addplot[black!60,thick,mark=x,mark size=1.6pt] table[x=lag,y=r]{figdata/research/05-stylised-facts-of-returns/acf.csv};
\addplot[omThm,very thick,dashed] table[x=lag,y=synth_abs]{figdata/research/05-stylised-facts-of-returns/acf.csv};
\legend{$|r|$ (US market),$r^2$ (US market),$r$ (US market),$|r|$ (synthetic market)}
\end{semilogxaxis}
\end{tikzpicture}

\omcaption{Autocorrelations of daily returns, absolute returns and squared returns of the US market, 1926--2026, and of the
absolute returns of the synthetic market factor. Returns are nearly uncorrelated; absolute returns stay correlated beyond a hundred
days, and more than squared returns (the \omterm{def:rs:stylised-facts-of-returns:taylor}{Taylor effect}). The synthetic factor's clustering fades after about fifty days. Data:
Kenneth R.~French data library (derived statistics); \texttt{firm.synthmkt}, seed 1.}\label{fig:rs:stylised-facts-of-returns:acf}
\end{omfigure}

\section{The leverage effect and gain--loss asymmetry}

\begin{definition}[Gain--loss asymmetry]\label{def:rs:stylised-facts-of-returns:gainloss}
\emph{Gain--loss asymmetry}\index{gain--loss asymmetry} is the observation that stock prices and indices have large falls but not
equally large rises: the distribution of returns is negatively skewed, and more so at horizons of weeks than of days.
\end{definition}

The century's daily skewness is $-0.16$; over a week it is $-0.76$ and over a month $-0.90$. The largest daily fall is the $-17.4\%$
of 19 October 1987; the largest daily rise, $+15.7\%$ on 15 March 1933, came in the most volatile market of the century. The leverage
effect (Book 4, chapter 18), the negative correlation between returns and later volatility, is visible in the same data: the
correlation between today's return and tomorrow's squared return is $-0.085$, against $+0.025$ between today's return and
yesterday's squared return. Falls raise volatility; volatility does not predict the sign of returns.

\section{Intraday seasonality}

\begin{definition}[Intraday seasonality, intraday volume profile]\label{def:rs:stylised-facts-of-returns:intraday}
\emph{Intraday seasonality}\index{intraday seasonality} is the regular dependence of volatility, volume, spreads and trading
activity on the time of day. The \emph{intraday volume profile}\index{intraday volume profile} is the average share of a day's
volume traded in each interval of the session.
\end{definition}

Equity markets trade most, and move most, just after the open and just before the close: Wood, McInish and Ord (1985) found
unusually high returns and standard deviations at the beginning and end of the day in minute-by-minute NYSE data, and Andersen and
Bollerslev (1997) model this intraday periodicity jointly with the persistence of volatility from day to day. The pattern matters to anyone who samples in clock time: a return over the first five minutes of the day and one over five
minutes at noon are draws from different distributions. The book's intraday simulator builds the profile in (chapter~2): on its day
the first quarter-hour trades 2.6 times the median quarter-hour's volume and the last 1.7 times. The chapter has no licensed intraday
data to measure it on real stocks; the profile of the simulator is an assumption, stated as one.

\section{Cross-sectional structure}

Stocks move together, and how much depends on the market's direction.

\begin{definition}[Exceedance correlation, correlation asymmetry]\label{def:rs:stylised-facts-of-returns:exceedance}
An \emph{exceedance correlation}\index{exceedance correlation} is a correlation between returns measured only on the days when a
conditioning variable (a market return, or the two returns themselves) lies beyond a threshold. \emph{Correlation
asymmetry}\index{correlation asymmetry} is the finding that exceedance correlations in falling markets exceed those in rising
markets at the same threshold.
\end{definition}

Longin and Solnik (2001) found, with extreme-value methods, that correlation between international equity markets increases in bear
markets but not in bull markets; Ang and Chen (2002) found the same asymmetry between US stock portfolios and the market. The
French data library's ten industry portfolios show it over the century (\cref{fig:rs:stylised-facts-of-returns:asym}): on days when
the market fell by more than one standard deviation their average pairwise correlation is 0.61; on days it rose as much, 0.56; the
gap is there at every threshold from zero to two standard deviations. Diversification across industries is worth least on the days
it is needed.

The second cross-sectional fact is a dominant factor. The correlation matrix of 200 simulated stocks over a year has a largest
eigenvalue of 44.5, 22\% of the total, against a Marchenko--Pastur upper edge of 3.58 for pure noise of the same shape (Book 4,
chapter 22); four eigenvalues lie above the edge: the market and the strongest industries.

\begin{omfigure}
\begin{tikzpicture}
\begin{axis}[width=10.5cm,height=5.2cm,xlabel={threshold (standard deviations of the market's daily return)},
  ylabel={average pairwise correlation},tick label style={font=\small},label style={font=\small},xmin=-0.05,xmax=2.05,ymin=0.1,
  ymax=0.7,grid=major,grid style={gray!25},legend columns=2,legend cell align=left,legend style={font=\footnotesize,at={(0.5,-0.32)},anchor=north,
  draw=none,fill=none,/tikz/every even column/.append style={column sep=8pt}},table/col sep=comma]
\addplot[omThm,very thick,mark=*] table[x=threshold,y=down]{figdata/research/05-stylised-facts-of-returns/asym.csv};
\addplot[omDef,very thick,mark=square*] table[x=threshold,y=up]{figdata/research/05-stylised-facts-of-returns/asym.csv};
\addplot[omThm!60,thick,dashed,mark=o] table[x=threshold,y=synth_down]{figdata/research/05-stylised-facts-of-returns/asym.csv};
\addplot[omDef!60,thick,dashed,mark=square] table[x=threshold,y=synth_up]{figdata/research/05-stylised-facts-of-returns/asym.csv};
\legend{10 US industries: down days,10 US industries: up days,synthetic stocks: down days,synthetic stocks: up days}
\end{axis}
\end{tikzpicture}

\omcaption{Average pairwise correlation on days when the market fell (down) or rose (up) by more than the threshold. The ten US
industry portfolios, 1926--2026, are more correlated on down days at every threshold (0.61 against 0.56 at one standard deviation).
The book's synthetic stocks show the opposite ordering: the simulator misses this fact. Data: Kenneth R.~French data library
(derived statistics); \texttt{firm.synthmkt}, seed 1.}\label{fig:rs:stylised-facts-of-returns:asym}
\end{omfigure}

\section{A scorecard for the book's simulator}

\texttt{firm.synthmkt}, the build of this chapter, generates the cross-section the rest of the book trades: a market factor with
GJR-GARCH volatility and Student-$t$ shocks, industry and style factors, stocks with heterogeneous specific volatility, earnings
jumps, and four planted sources of expected return. The table scores it against the century of US data on the facts of this chapter.

\begin{center}
\small
\begin{tabular}{lcccc}
\toprule
& real & \multicolumn{3}{c}{synthetic (\texttt{firm.synthmkt})} \\
\cmidrule(lr){2-2}\cmidrule(lr){3-5}
fact (daily returns) & US market & market & EW index & median stock \\
\midrule
kurtosis & 19.1 & 20.5 & 18.0 & 8.7 \\
left tail index (Hill, largest 2\%) & 2.90 & 2.83 & 2.88 & 3.35 \\
autocorrelation of $|r|$, lag 1 & 0.30 & 0.23 & 0.25 & 0.10 \\
autocorrelation of $|r|$, lag 20 & 0.22 & 0.21 & 0.22 & 0.08 \\
autocorrelation of $|r|$, lag 100 & 0.13 & $-0.02$ & $-0.03$ & $-0.01$ \\
autocorrelation of $r$, lag 1 & 0.05 & $-0.01$ & 0.01 & $-0.03$ \\
correlation of $r_t$ with $r^2_{t+1}$ & $-0.085$ & $-0.111$ & $-0.103$ & $-0.024$ \\
kurtosis of monthly log returns & 9.5 & 3.8 & 3.9 & 3.5 \\
\bottomrule
\end{tabular}
\end{center}

\begin{remark}[What the simulator can and cannot be used for]\label{rem:rs:stylised-facts-of-returns:scorecard}
The simulator reproduces the daily tails, the clustering over a month and the leverage effect, which is what the backtests,
predictor measurements and risk models of this book need. It misses three facts: volatility memory beyond about fifty days (GARCH
decays exponentially), the heavy tails of monthly returns (it has no crashes longer than a day), and \omterm{def:rs:stylised-facts-of-returns:exceedance}{correlation asymmetry}. Research
whose answer depends on them (tail hedging, drawdown risk over months, the value of diversification in a sell-off) must not be
validated on it.
\end{remark}

\section{Tutorial: scoring a simulator}\label{tut:rs:stylised-facts-of-returns:scorecard}

\begin{tutorial}
\textbf{Goal.} Compute the \omterm{def:rs:stylised-facts-of-returns:fact}{stylised facts} of a return series, apply them to the century of US data and to the synthetic market,
and fill in the scorecard. \textbf{End state:} the table above and
\cref{fig:rs:stylised-facts-of-returns:agg,fig:rs:stylised-facts-of-returns:acf,fig:rs:stylised-facts-of-returns:asym}.
\begin{tutsteps}
\item \textbf{The statistics.} One function computes every row of the scorecard for one series; the Hill estimator is Book 4's.
\omcode{code/research/05-stylised-facts-of-returns/python/rs_stylised.py}{55}{64}{The scorecard of one return series.}\label{lst:rs:stylised-facts-of-returns:facts}
\item \textbf{The market factor.} GJR-GARCH(1,1): negative shocks raise tomorrow's variance more than positive ones; the
Student-$t$ shocks give the tails that clustering alone does not.
\omcode{code/firm/synthmkt/firm_synthmkt.py}{156}{167}{The synthetic market factor.}\label{lst:rs:stylised-facts-of-returns:garch}
\item \textbf{Run} \texttt{synth\_series()} for the synthetic rows, \texttt{french\_summary()} for the real one, and
\texttt{fig\_stylised.py}. The real statistics come from \texttt{rs\_fetch\_french.py}, which keeps only derived statistics of the
library's data.
\end{tutsteps}
\textbf{What to change next.} Raise \texttt{garch\_beta} to 0.96 and \texttt{garch\_alpha} to 0.01 (same persistence, slower
decay) and see the autocorrelation at lag 100 move; make betas larger on down days and look for \omterm{def:rs:stylised-facts-of-returns:exceedance}{correlation asymmetry}, and for what
it does to the mean return.
\end{tutorial}

\section{Build: the synthetic equity market}\label{bld:rs:stylised-facts-of-returns:synthmkt}

\begin{build}
\bfield{Purpose} The book's cross-section with the truth attached. Chapters 6--15 measure predictors on it, Part~IV backtests on it,
Part~V builds risk models and portfolios on it, and Books 8 and 9 run their equity strategy tutorials on it.
\bfield{Interface} \texttt{MarketConfig} (size, length, seed; market, industry and style factor parameters; specific volatility; the
four planted alpha strengths; delisting, merger, split, dividend and filing rules); \texttt{simulate(cfg)} returning a \texttt{Panel}:
daily \texttt{ret}, \texttt{price}, \texttt{volume}, \texttt{shares}, \texttt{listed}; \texttt{industry}, \texttt{beta},
\texttt{style\_x}, \texttt{spec\_vol}; the factor paths and the market's conditional variance; \texttt{alpha} (the planted expected
return for the next day, by component); delistings, splits, dividends, earnings with their true surprises, point-in-time fundamentals
with \omterm{def:rs:point-in-time-data-and-the-biases:vintage}{restatements}; a \texttt{secmaster}.
\bfield{Rules} Deterministic for a seed. The planted expected return for day $t+1$ is recorded at day $t$, so that a signal known at
the close of $t$ can be scored against it. A delisted name is replaced the next day. Fundamentals are filed after the quarter's end
and \omterm{def:rs:point-in-time-data-and-the-biases:vintage}{restatements} come later still.
\bfield{Acceptance tests} \texttt{code/firm/synthmkt/tests/}: determinism and bookkeeping (listed count, \omterm{def:rs:universe-symbology-and-corporate-actions:delisting}{delisting returns}, no return
after a delisting); agreement with the \omterm{def:rs:universe-symbology-and-corporate-actions:master}{security master}; point-in-time fundamentals through \texttt{firm.pit}; over three seeds, market
kurtosis above 5, positive volatility clustering at twenty days, a negative leverage correlation and positive reversal and momentum
information coefficients; a dominant first eigenvalue.
\bfield{Stretch} Downside betas with a compensating drift (\omterm{def:rs:stylised-facts-of-returns:exceedance}{correlation asymmetry}); a slowly varying volatility component
(long memory); multi-day crashes.
\end{build}

\begin{omsources}
\item R.~Cont, ``Empirical properties of asset returns: stylized facts and statistical issues'', \emph{Quantitative Finance} 1(2),
2001.
\item P.~Gopikrishnan, M.~Meyer, L.~A.~N.~Amaral and H.~E.~Stanley, ``Inverse cubic law for the distribution of stock price
variations'', \emph{European Physical Journal B} 3, 1998.
\item C.~W.~J.~Granger and Z.~Ding, ``Varieties of long memory models'', \emph{Journal of Econometrics} 73, 1996; S.~J.~Taylor,
\emph{Modelling Financial Time Series}, Wiley, 1986.
\item R.~A.~Wood, T.~H.~McInish and J.~K.~Ord, ``An investigation of transactions data for NYSE stocks'', \emph{Journal of Finance}
40(3), 1985; T.~G.~Andersen and T.~Bollerslev, ``Intraday periodicity and volatility persistence in financial markets'',
\emph{Journal of Empirical Finance} 4, 1997.
\item F.~Longin and B.~Solnik, ``Extreme correlation of international equity markets'', \emph{Journal of Finance} 56(2), 2001;
A.~Ang and J.~Chen, ``Asymmetric correlations of equity portfolios'', \emph{Journal of Financial Economics} 63(3), 2002.
\item B.~Mandelbrot, ``The variation of certain speculative prices'', \emph{Journal of Business} 36(4), 1963.
\item Kenneth R.~French, Data Library: Fama/French 3 factors and 10 industry portfolios, daily.
\item Federal Reserve History, ``Stock Market Crash of 1987''.
\end{omsources}

\section{Exercises}

\begin{exercise}[$\star$]\label{exo:rs:stylised-facts-of-returns:1}
A normal law with the century's standard deviation of 1.08\% a day: what is the probability of a day below $-17.4\%$, and how many
such days would you expect in 26\,296?
\end{exercise}

\begin{exercise}[$\star$]\label{exo:rs:stylised-facts-of-returns:2}
If daily returns were independent with kurtosis 20, what would the kurtosis of their sum over 5, 21 and 63 days be? Compare with the
century's 11.5, 9.5 and 7.3.
\end{exercise}

\begin{exercise}[$\star$]\label{exo:rs:stylised-facts-of-returns:3}
A power-law tail with index 3: if one day in a thousand falls by more than 4\%, how often does a day fall by more than 8\%? And with an
index of 4?
\end{exercise}

\begin{exercise}[$\star\star$]\label{exo:rs:stylised-facts-of-returns:4}
Show that for a Student-$t$ law with $\nu > 4$ degrees of freedom, scaled to unit variance, the kurtosis is $3 + 6/(\nu - 4)$. What $\nu$
matches a kurtosis of 19?
\end{exercise}

\begin{exercise}[$\star\star$]\label{exo:rs:stylised-facts-of-returns:5}
The ten industries have an average pairwise correlation of 0.61 on the market's down days and 0.56 on its up days (one standard
deviation). With equal weights and equal volatilities $\sigma$, what is the volatility of the equal-weighted portfolio relative to
$\sigma$ in each case?
\end{exercise}

\begin{exercise}[$\star\star$]\label{exo:rs:stylised-facts-of-returns:6}
The GJR-GARCH market factor has $\alpha = 0.02$, $\gamma = 0.10$, $\beta = 0.92$. What is its persistence, and the half-life of a
variance shock? Why does its autocorrelation of absolute returns fade after about fifty days?
\end{exercise}

\begin{exercise}[$\star\star\star$]\label{exo:rs:stylised-facts-of-returns:7}
\emph{Coding.} Simulate the synthetic market with \texttt{garch\_alpha=0.01}, \texttt{garch\_gamma=0.05},
\texttt{garch\_beta=0.965} (the same persistence) and compare the autocorrelations of absolute market returns at lags 20 and 100
with the default's.
\end{exercise}

\begin{exercise}[$\star\star\star$]\label{exo:rs:stylised-facts-of-returns:8}
\emph{Find the flaw.} ``We validated our tail-hedging strategy on a simulated market calibrated to the daily kurtosis and volatility
clustering of the S\&P~500. It pays for itself over any ten-year window.''
\end{exercise}

\section{Problem: Is the Simulator Realistic?}

\begin{problem}[{Weekend problem --- a scorecard of the book's synthetic market against a century of data}]\label{pb:rs:stylised-facts-of-returns:1}
The century of daily US market excess returns (French data library, derived statistics) and \texttt{firm.synthmkt} with its default
configuration (1\,000 names, ten years, seed 1).

\medskip\noindent\textbf{Part I --- Tails.}
\begin{enumerate}
\item What are the daily kurtoses of the US market and of the synthetic market factor?
\item What left tail indices does the Hill estimator give each on its largest 2\% of losses?
\item How many days of the century lie beyond five standard deviations, and how many would a normal law predict?
\item How many standard deviations of the previous five years was 19 October 1987?
\item What kurtosis does the synthetic median stock have, and why is it lower than the index's?
\end{enumerate}

\medskip\noindent\textbf{Part II --- Dependence.}
\begin{enumerate}[resume]
\item What are the autocorrelations of absolute returns at lags 1, 20 and 100, real and synthetic?
\item Where does the synthetic market fail, and why?
\item Does the synthetic market show the leverage effect? Compare the correlations of $r_t$ with $r^2_{t+1}$.
\item What are the kurtoses of monthly log returns, and what explains the gap?
\item Does the century show the \omterm{def:rs:stylised-facts-of-returns:taylor}{Taylor effect} at lag 20?
\end{enumerate}

\medskip\noindent\textbf{Part III --- The cross-section.}
\begin{enumerate}[resume]
\item What are the ten industries' average correlations on down and up days at one standard deviation?
\item What does the synthetic market give, and what does that mean?
\item What is the largest eigenvalue of 200 synthetic stocks' correlation over a year, its share, and the Marchenko--Pastur edge?
\item How many eigenvalues lie above the edge, and what are they?
\end{enumerate}

\medskip\noindent\textbf{Part IV --- Verdict.}
\begin{enumerate}[resume]
\item Which research of this book may use the simulator without reservation?
\item Which research must not be validated on it?
\item What single change would most improve it for a tail-risk study?
\item State the \emph{named result}: the scorecard's three rows the simulator matches best and the three it misses.
\item Why does a simulator used for research need a scorecard at all?
\item In one sentence: what is a \omterm{def:rs:stylised-facts-of-returns:fact}{stylised fact} good for?
\end{enumerate}
\end{problem}

\section{Interview questions}

\begin{interviewq}[$\star$ \iqroles{researcher, trader}]\label{iq:rs:stylised-facts-of-returns:1}
Name four \omterm{def:rs:stylised-facts-of-returns:fact}{stylised facts} of equity returns.
\end{interviewq}

\begin{interviewq}[$\star\star$ \iqroles{researcher, risk}]\label{iq:rs:stylised-facts-of-returns:2}
Daily returns are nearly uncorrelated. Does that make them independent? What would you check?
\end{interviewq}

\begin{interviewq}[$\star\star$ \iqroles{risk, researcher}]\label{iq:rs:stylised-facts-of-returns:3}
Why does diversification across sectors disappoint in a crash?
\end{interviewq}

\begin{interviewq}[$\star\star$ \iqroles{researcher, mle}]\label{iq:rs:stylised-facts-of-returns:4}
You want a simulator to test a trading strategy. How do you decide whether it is realistic enough?
\end{interviewq}

\begin{interviewq}[$\star\star\star$ \iqroles{researcher}]\label{iq:rs:stylised-facts-of-returns:5}
If daily returns have kurtosis 20 and are independent, what is the kurtosis of monthly returns? The data say about 9. What does that
tell you?
\end{interviewq}

\begin{interviewq}[$\star\star$ \iqroles{trader, researcher}]\label{iq:rs:stylised-facts-of-returns:6}
What is the leverage effect, and how would you measure it?
\end{interviewq}
